\section{A toy model}
\label{sec:toy}

The toy model isolates the interaction between population growth, market power in
current AI services, and the automation of AI research. Time is continuous. There is
one final good, a representative dynasty, competitive final-good firms, and one
integrated frontier developer. The developer monopolistically supplies AI services
and invests in the next generation of AI. The initial economy already uses AI
services in production; the model does not describe their first commercial adoption.

\subsection{Household}

Population grows at the exogenous rate $n$:
\begin{equation}
    \dot N_t=nN_t,
    \qquad N_0>0,\quad n\geq0.
    \label{eq:population}
\end{equation}
The representative dynasty has preferences
\begin{equation}
    \int_0^\infty e^{-\rho t}N_t\log c_t\dd t,
    \qquad
    c_t\equiv\frac{C_t}{N_t},
    \qquad \rho>n,
    \label{eq:utility}
\end{equation}
where $C_t$ is aggregate consumption. Each household member supplies one unit of
labor inelastically. Labor is allocated between final production, $L_{Y,t}$, and
human AI research, $H_t$:
\begin{equation}
    L_{Y,t}+H_t=N_t.
    \label{eq:labor-clearing}
\end{equation}
The dynasty owns the capital stock and the frontier developer. Capital evolves as
\begin{equation}
    \dot K_t=I_t-\delta_KK_t.
    \label{eq:capital-law}
\end{equation}
Let $k_t\equiv K_t/N_t$. The per-capita budget constraint is
\begin{equation}
    \dot k_t
    =
    (r_t-\delta_K-n)k_t+w_t
    +\frac{\pi^X_t-w_tH_t-\xi M_t}{N_t}
    -c_t,
    \label{eq:household-budget}
\end{equation}
where $\pi^X_t$ is the developer's operating profit before research expenditure.
The household's Euler equation is
\begin{equation}
    \frac{\dot c_t}{c_t}=r_t-\delta_K-\rho.
    \label{eq:euler}
\end{equation}

\subsection{Final-good firms}

Competitive final-good firms combine physical capital with a composite of production
labor and AI services:
\begin{align}
    Y_t
    &=
    K_t^\alpha Z_t^{1-\alpha},
    \qquad \alpha\in(0,1),
    \label{eq:final-production}\\
    Z_t
    &=
    \left[
        (1-\omega)L_{Y,t}^{(\sigma-1)/\sigma}
        +\omega X_t^{(\sigma-1)/\sigma}
    \right]^{\sigma/(\sigma-1)},
    \qquad
    \omega\in(0,1),\quad \sigma>0.
    \label{eq:service-composite}
\end{align}
The elasticity of substitution between capital and the service composite is one.
The elasticity between production labor and AI services is $\sigma$. The
Cobb--Douglas composite is obtained by continuity when $\sigma=1$.

Define the contribution of AI services to the CES composite as
\begin{equation}
    s^X_t
    \equiv
    \frac{\omega X_t^{(\sigma-1)/\sigma}}
    {(1-\omega)L_{Y,t}^{(\sigma-1)/\sigma}
        +\omega X_t^{(\sigma-1)/\sigma}}.
    \label{eq:ai-composite-share}
\end{equation}
Let $p^X_t$ denote the price of AI services. Profit maximization gives
\begin{align}
    r_t
    &=\alpha\frac{Y_t}{K_t},
    \label{eq:capital-demand}\\
    w_t
    &=(1-\alpha)(1-s^X_t)\frac{Y_t}{L_{Y,t}},
    \label{eq:labor-demand}\\
    p^X_t
    &=(1-\alpha)s^X_t\frac{Y_t}{X_t},
    \label{eq:ai-demand}
\end{align}
and relative demand satisfies
\begin{equation}
    \frac{X_t}{L_{Y,t}}
    =
    \left[
        \frac{\omega}{1-\omega}
        \frac{w_t}{p^X_t}
    \right]^\sigma.
    \label{eq:relative-ai-demand}
\end{equation}

\subsection{Integrated AI monopolist}

The variable $U_t$ is a flow of \emph{inference compute}: standardized
computational resources used during period $t$ to run an existing model. It can be
measured in FLOPs, GPU-hours, or equivalent units of a reference processor per unit
of time; it is not a stock of chips. Let $\xi>0$ denote the common unit resource cost
of standardized compute. It summarizes the capital, energy, and other resources
needed to supply one unit of compute, independently of whether that unit is used for
inference or research.

The variable $X_t$ is a flow of quality-adjusted AI services used by final-good
firms. Its empirical counterpart could be benchmark-equivalent tokens, predictions
adjusted for accuracy, or equivalent task hours. This distinction follows the
measurement approach in \citet{korinekmckelvey2026}, who combine inference compute
with model capital and adjust inference output for quality, and is consistent with
the intelligence tiers documented by \citet{demireretal2025}.

The stock $A_t$ measures frontier AI capability. Under the normalization used here,
it also measures quality-adjusted services produced per unit of inference compute:
\begin{equation}
    X_t=A_tU_t.
    \label{eq:ai-services}
\end{equation}
Equation~\eqref{eq:ai-services} is an efficiency-units specification, not a
technological identity. Conditional on $A_t$, it imposes constant returns to
inference compute and compresses task coverage, quality, speed, and inference
efficiency into one scalar.

\begin{remark}[Capability and inference efficiency]
The same state variable $A_t$ governs broad frontier capability and the productivity
of inference compute. This restriction keeps the service-market and research
mechanisms connected and makes the benchmark tractable, but it need not hold
quantitatively. A natural extension is $X_t=B(A_t)U_t$, with a separate mapping from
capability to inference efficiency. We retain $B(A)=A$ in the baseline and keep this
distinction as a robustness extension.
\end{remark}

The integrated developer recognizes the inverse demand generated by the final-good
sector and solves
\begin{equation}
    \max_{X_t\geq0}
    \left\{
        p^X_t(X_t;K_t,L_{Y,t})X_t
        -\xi\frac{X_t}{A_t}
    \right\}.
    \label{eq:monopoly-service-problem}
\end{equation}
Its operating profit is
\begin{equation}
    \pi^X_t=p^X_tX_t-\xi\frac{X_t}{A_t}.
    \label{eq:operating-profit}
\end{equation}

\subsection{Frontier research}

New AI capability is produced with effective research services, $E_t$:
\begin{align}
    \dot A_t
    &=
    \chi A_t^\phi E_t^\eta,
    \qquad \chi>0,\quad \phi<1,\quad \eta\in(0,1),
    \label{eq:ideas}\\
    E_t
    &=
    \left[
        (1-\nu)H_t^{(\theta-1)/\theta}
        +\nu M_t^{(\theta-1)/\theta}
    \right]^{\theta/(\theta-1)},
    \qquad
    \nu\in(0,1),\quad \theta>1.
    \label{eq:effective-research}
\end{align}
Here, $M_t$ is a flow of standardized automated-research services. It includes
training, post-training, experiments, synthetic-data generation, and inference
performed by AI agents engaged in R\&D. The baseline normalizes one unit of these
services to one unit in the research aggregator and assigns it the common resource
cost $\xi$. A capability-dependent conversion from compute into automated-research
services is left for an extension.

The elasticity $\theta$ governs substitution between human and automated research.
The restriction $\theta>1$ makes the two inputs gross substitutes and allows
positive research output when $H_t=0$. It does not impose that human research
disappears. That outcome depends on the equilibrium path of the wage relative to
$\xi$. The exponent $\phi$ captures how the existing stock of knowledge affects the
productivity of both research inputs.

\begin{remark}[Capability-dependent research productivity]
An extension replaces $M_t$ inside \eqref{eq:effective-research} with
$\lambda(A_t)M_t$. Its effective unit cost is then $\xi/\lambda(A_t)$, and the
automation condition depends on $w_t\lambda(A_t)/\xi$ rather than $w_t/\xi$. In the
AI-dominated limit, a power function $\lambda(A)=\lambda_0A^\psi$ changes the
capability exponent from $\phi$ to $\phi+\psi\eta$. The baseline sets
$\lambda(A)=1$ to isolate substitution between human and automated research from
capability-dependent improvements in research efficiency.
\end{remark}

Let $q_t$ be the developer's private shadow value of capability and let
$F(A_t,H_t,M_t)$ denote the right-hand side of \eqref{eq:ideas}. Its current-value
Hamiltonian is
\begin{equation}
    \mathcal H^P_t
    =
    \pi^X_t-w_tH_t-\xi M_t
    +q_tF(A_t,H_t,M_t).
    \label{eq:private-hamiltonian}
\end{equation}

\subsection{Social planner}

Let $Q_t$ denote the social shadow value of capability. Unlike $q_t$, $Q_t$ includes
consumer surplus from a larger supply of AI services and knowledge benefits that the
developer cannot appropriate. The comparison between $Q_t$ and $q_t$ isolates the
knowledge and appropriability wedge.

\subsection{Equilibrium}

A decentralized equilibrium is a sequence
\[
    \{C_t,c_t,I_t,K_t,N_t,L_{Y,t},H_t,X_t,U_t,M_t,A_t\}_{t\geq0}
\]
and prices
\[
    \{r_t,w_t,p^X_t,q_t\}_{t\geq0}
\]
such that:
\begin{enumerate}[label=(\roman*)]
    \item the representative dynasty maximizes \eqref{eq:utility}, taking prices and
    developer profits as given;
    \item final-good firms maximize profits under
    \eqref{eq:final-production}--\eqref{eq:service-composite};
    \item the integrated developer chooses current AI services and frontier research
    to maximize its market value;
    \item population, capital, and capability follow
    \eqref{eq:population}, \eqref{eq:capital-law}, and \eqref{eq:ideas};
    \item labor satisfies \eqref{eq:labor-clearing}; and
    \item the aggregate resource constraint is
    \begin{equation}
        C_t+I_t+\xi(U_t+M_t)=Y_t.
        \label{eq:resource-constraint}
    \end{equation}
\end{enumerate}
There is no common physical-capacity constraint between $U_t$ and $M_t$. The two
uses of compute compete through the final-good resource constraint, not through an
additional exogenous cap.

\subsection{Solution}

\subsubsection{Wages and the labor share under an AI expansion}

Let
\begin{equation}
    \ell^Y_t\equiv\frac{w_tL_{Y,t}}{Y_t}
    =(1-\alpha)(1-s^X_t)
    \label{eq:production-labor-share}
\end{equation}
denote production labor's share of final output. The wage and this share need not
move in the same direction because the wage depends on both the production-labor
share and output per production worker.

\begin{proposition}[The wage and production-labor-share thresholds]
\label{prop:wage-share-thresholds}
Holding $K_t$ and $L_{Y,t}$ fixed, an expansion of AI services satisfies
\begin{align}
    \frac{\partial\log Y_t}{\partial\log X_t}
    &=(1-\alpha)s^X_t>0,
    \label{eq:output-ai-elasticity}\\
    \frac{\partial\log \ell^Y_t}{\partial\log X_t}
    &=-\left(1-\frac{1}{\sigma}\right)s^X_t,
    \label{eq:labor-share-ai-elasticity}\\
    \frac{\partial\log w_t}{\partial\log X_t}
    &=\left(\frac{1}{\sigma}-\alpha\right)s^X_t.
    \label{eq:wage-ai-elasticity}
\end{align}
Consequently, the production-labor share falls if and only if $\sigma>1$, whereas
the wage rises if and only if $\sigma<1/\alpha$. Since $\alpha\in(0,1)$, there is
a nonempty interval
\begin{equation}
    1<\sigma<\frac{1}{\alpha}
    \label{eq:wage-share-decoupling-region}
\end{equation}
in which an expansion of AI raises the wage while reducing production labor's share
of output.
\end{proposition}

\begin{proof}
The CES expenditure share in \eqref{eq:ai-composite-share} implies
\[
    \frac{\partial\log(1-s^X_t)}{\partial\log X_t}
    =
    -\left(1-\frac{1}{\sigma}\right)s^X_t.
\]
The output elasticity follows directly from
\eqref{eq:final-production}--\eqref{eq:service-composite}. Combining these two
expressions with \eqref{eq:labor-demand} yields
\eqref{eq:labor-share-ai-elasticity} and \eqref{eq:wage-ai-elasticity}.
\end{proof}

Because researchers and production workers receive the same wage, the aggregate
labor-income share is $w_tN_t/Y_t=(N_t/L_{Y,t})\ell^Y_t$. Conditional on a fixed
allocation of labor between production and research, its elasticity with respect to
$X_t$ is also given by \eqref{eq:labor-share-ai-elasticity}. Along an equilibrium
transition, however, changes in $H_t/N_t$ provide an additional margin. The
production-block threshold should therefore not be read as a complete
general-equilibrium result for the aggregate labor share.

Equation~\eqref{eq:wage-ai-elasticity} decomposes the wage response into a
productivity effect and a displacement effect:
\begin{equation}
    \frac{\partial\log w_t}{\partial\log X_t}
    =
    \underbrace{(1-\alpha)s^X_t}_{\text{productivity effect}}
    -
    \underbrace{\left(1-\frac{1}{\sigma}\right)s^X_t}
    _{\text{displacement effect}}.
    \label{eq:wage-decomposition}
\end{equation}
The first term is the increase in output per production worker. The second is the
decline in labor's income share. This decomposition maps the nested production
function into the productivity and displacement channels emphasized by
\citet{acemoglurestrepo2019}.

The two thresholds are not unique to this paper. \citet{yango2026} derives the same
conditions in a nested technology with physical and AI capital, while
\citet{autorkausik2026} establish the broader possibility of falling labor shares
and rising wages under constant returns to scale. Proposition
\ref{prop:wage-share-thresholds} serves here as a benchmark for the model's distinct
dynamic and market-structure mechanisms.

Allowing all quantities to change, the instantaneous wage growth rate obeys the
identity
\begin{equation}
    g_{w,t}
    =
    \alpha\left(g_{K,t}-g_{L_Y,t}\right)
    +
    s^X_t\left(\frac{1}{\sigma}-\alpha\right)
    \left(g_{X,t}-g_{L_Y,t}\right).
    \label{eq:dynamic-wage-decomposition}
\end{equation}
When the research-employment share is constant, $g_{L_Y,t}=n$. Capital deepening can
therefore raise wages even when the direct expansion of AI services lowers them.
Conversely, $\sigma<1/\alpha$ is not sufficient for wage growth along an equilibrium
transition if capital per production worker contracts. The monopoly problem below
determines the supply of $X_t$, while household saving and the resource constraint
determine $K_t$; the equilibrium wage response depends on both paths.

\subsubsection{Monopoly supply of AI services}

Holding $K_t$ and $L_{Y,t}$ fixed, the absolute elasticity of inverse demand faced by
the monopolist is
\begin{equation}
    \varepsilon^X_t
    \equiv
    -\frac{\partial\log p^X_t}{\partial\log X_t}
    =
    \frac{1-s^X_t}{\sigma}+\alpha s^X_t.
    \label{eq:inverse-demand-elasticity}
\end{equation}
For an interior solution with $\varepsilon^X_t<1$, the first-order condition is
\begin{equation}
    p^X_t(1-\varepsilon^X_t)=\frac{\xi}{A_t}.
    \label{eq:monopoly-ai-price}
\end{equation}
The markup and operating rent are
\begin{align}
    \mu^X_t
    &\equiv
    \frac{p^X_t}{\xi/A_t}
    =
    \frac{1}{1-\varepsilon^X_t}>1,
    \label{eq:monopoly-markup}\\
    \pi^X_t
    &=
    \varepsilon^X_t p^X_tX_t.
    \label{eq:monopoly-rents}
\end{align}
Conditional on $A_t$, $K_t$, and $L_{Y,t}$, the monopolist supplies less AI than an
allocation that prices services at marginal cost. The household receives the rents,
so the toy model retains the output distortion and the innovation incentive from
market power while abstracting from concentrated ownership.

\subsubsection{Research input choice and gradual automation}

The marginal products of the two research inputs in
\eqref{eq:effective-research} are
\begin{equation}
    E_{H,t}
    =(1-\nu)\left(\frac{E_t}{H_t}\right)^{1/\theta},
    \qquad
    E_{M,t}
    =\nu\left(\frac{E_t}{M_t}\right)^{1/\theta}.
    \label{eq:research-ces-marginal-products}
\end{equation}
The developer's first-order conditions are therefore
\begin{align}
    w_t
    &=
    q_t\chi\eta A_t^\phi E_t^{\eta-1}
    (1-\nu)\left(\frac{E_t}{H_t}\right)^{1/\theta},
    \label{eq:foc-human}\\
    \xi
    &=
    q_t\chi\eta A_t^\phi E_t^{\eta-1}
    \nu\left(\frac{E_t}{M_t}\right)^{1/\theta}.
    \label{eq:private-compute-foc}
\end{align}
Their ratio determines the research-input mix:
\begin{equation}
    \frac{M_t}{H_t}
    =
    \left[
        \frac{\nu}{1-\nu}\frac{w_t}{\xi}
    \right]^\theta.
    \label{eq:research-input-ratio}
\end{equation}

The unit cost of effective research services is
\begin{equation}
    P^E_t
    =
    \left[
        (1-\nu)^\theta w_t^{1-\theta}
        +\nu^\theta\xi^{1-\theta}
    \right]^{1/(1-\theta)}.
    \label{eq:research-unit-cost}
\end{equation}
Conditional input demands are
\begin{equation}
    H_t=(1-\nu)^\theta
    \left(\frac{P^E_t}{w_t}\right)^\theta E_t,
    \qquad
    M_t=\nu^\theta
    \left(\frac{P^E_t}{\xi}\right)^\theta E_t.
    \label{eq:research-conditional-demands}
\end{equation}
The developer can consequently choose the scale of research in two stages. For
$\eta\in(0,1)$, its optimal scale satisfies
\begin{equation}
    P^E_t=q_t\chi\eta A_t^\phi E_t^{\eta-1},
    \qquad
    E_t=
    \left(
        \frac{q_t\chi\eta A_t^\phi}{P^E_t}
    \right)^{1/(1-\eta)}.
    \label{eq:research-scale}
\end{equation}

Define the contribution of automated services to effective research as
\begin{equation}
    s^M_t
    \equiv
    \frac{\nu M_t^{(\theta-1)/\theta}}
    {(1-\nu)H_t^{(\theta-1)/\theta}
        +\nu M_t^{(\theta-1)/\theta}}.
    \label{eq:automated-research-share}
\end{equation}
At an interior cost-minimizing allocation, this is also the share of research
expenditure devoted to automated services:
\begin{equation}
    s^M_t=\frac{\xi M_t}{w_tH_t+\xi M_t}.
    \label{eq:automated-research-cost-share}
\end{equation}
Equations \eqref{eq:research-input-ratio} and
\eqref{eq:automated-research-share} imply
\begin{equation}
    \frac{s^M_t}{1-s^M_t}
    =
    \left(\frac{\nu}{1-\nu}\right)^\theta
    \left(\frac{w_t}{\xi}\right)^{\theta-1}.
    \label{eq:automation-odds}
\end{equation}
More generally, if the unit cost of automated research varies at rate $g_{\xi,t}$,
the automation share evolves according to
\begin{equation}
    \dot s^M_t
    =
    s^M_t(1-s^M_t)(\theta-1)
    \left(g_{w,t}-g_{\xi,t}\right).
    \label{eq:automation-share-dynamics}
\end{equation}
The baseline holds $\xi$ constant, so $g_{\xi,t}=0$.

\begin{proposition}[Gradual research automation]
\label{prop:gradual-automation}
Suppose $\theta>1$. If $w_t/\xi\to\infty$, then $s^M_t\to1$ and
$H_t/E_t\to0$. If $w_t/\xi$ remains bounded, $\theta>1$ alone does not imply that
the human contribution disappears.
\end{proposition}

Proposition~\ref{prop:gradual-automation} replaces a discrete automation threshold
with a continuous equilibrium transition. The research elasticity determines
whether human research can become asymptotically inessential. The production
elasticity $\sigma$ affects whether that limit is reached by determining the wage
path. The private shadow value continues to satisfy
\begin{equation}
    (r_t-\delta_K)q_t
    =
    \pi^X_{A,t}+q_tF_{A,t}+\dot q_t,
    \qquad
    F_{A,t}=\phi\frac{\dot A_t}{A_t}.
    \label{eq:private-costate}
\end{equation}

\subsubsection{Human-dominated research}

When $s^M_t$ is close to zero,
\begin{equation}
    E_t
    \simeq
    (1-\nu)^{\theta/(\theta-1)}H_t.
    \label{eq:human-dominated-research}
\end{equation}
If a constant share $s_H\in(0,1)$ of the population works in research,
\begin{equation}
    H_t=s_HN_t,
    \qquad
    L_{Y,t}=(1-s_H)N_t,
    \label{eq:research-share}
\end{equation}
the ideas equation implies
\begin{equation}
    g_A^H
    =
    \frac{\eta n}{1-\phi}.
    \label{eq:human-dominated-ga}
\end{equation}
This is a semi-endogenous benchmark for the early, human-dominated phase. It is not
generally a permanent balanced-growth path once automated research is feasible. If
the wage grows relative to $\xi$, equation \eqref{eq:automation-share-dynamics}
moves the economy away from this phase.

\subsubsection{AI-dominated research dynamics}

When $s^M_t\to1$, the research aggregator satisfies
\begin{equation}
    E_t
    \simeq
    \nu^{\theta/(\theta-1)}M_t.
    \label{eq:ai-dominated-research}
\end{equation}
The asymptotic law of motion is therefore
\begin{equation}
    \dot A_t
    =
    \kappa A_t^\phi M_t^\eta,
    \qquad
    \kappa\equiv
    \chi\nu^{\eta\theta/(\theta-1)}.
    \label{eq:ai-dominated-ideas}
\end{equation}
Its growth rate obeys
\begin{equation}
    \frac{\dot g_{A,t}}{g_{A,t}}
    =
    (\phi-1)g_{A,t}+\eta g_{M,t}.
    \label{eq:ai-growth-rate-dynamics}
\end{equation}
If $g_{M,t}$ converges to a positive constant and $\phi<1$, a constant capability
growth rate must satisfy
\begin{equation}
    g_A^\infty
    =
    \frac{\eta g_M^\infty}{1-\phi}.
    \label{eq:ai-dominated-ga}
\end{equation}
Research growth is no longer tied directly to population growth. It is tied to the
growth of automated-research services, whose equilibrium path depends on output,
profits, and the private value of capability.

\subsubsection{A closed-form AI-dominated path under Cobb--Douglas production}

The Cobb--Douglas limit of the production CES provides a closed-form benchmark.
Define
\begin{equation}
    \beta\equiv(1-\alpha)\omega,
    \qquad
    \upsilon\equiv\frac{\beta}{1-\alpha-\beta}.
    \label{eq:beta-definition}
\end{equation}
Production becomes
\begin{equation}
    Y_t=K_t^\alpha L_{Y,t}^{1-\alpha-\beta}X_t^\beta.
    \label{eq:cd-production}
\end{equation}
The monopolist chooses
\begin{align}
    p^X_t&=\frac{\xi}{\beta A_t},
    &
    U_t&=\frac{\beta^2}{\xi}Y_t,
    &
    X_t&=\frac{\beta^2A_t}{\xi}Y_t,
    \label{eq:cd-monopoly}\\
    \pi^X_t&=\beta(1-\beta)Y_t.
    \label{eq:cd-monopoly-rents}
\end{align}
Substitution into production gives
\begin{equation}
    Y_t
    =
    \left(\frac{\beta^2}{\xi}\right)^{\beta/(1-\beta)}
    A_t^{\beta/(1-\beta)}
    K_t^{\alpha/(1-\beta)}
    L_{Y,t}^{(1-\alpha-\beta)/(1-\beta)}.
    \label{eq:reduced-production}
\end{equation}

Let $\gamma^\infty$ denote the common growth rate of output, capital,
consumption, and inference compute per capita. Along an AI-dominated path with a
constant research-resource share $m_R\equiv\xi M_t/Y_t$,
\begin{equation}
    \gamma^\infty=\upsilon g_A^\infty,
    \qquad
    g_M^\infty=n+\gamma^\infty.
    \label{eq:cd-ai-growth-relations}
\end{equation}
Combining \eqref{eq:ai-dominated-ga} and
\eqref{eq:cd-ai-growth-relations} yields
\begin{align}
    g_A^\infty
    &=
    \frac{\eta n}
    {1-\phi-\eta\upsilon},
    \label{eq:cd-ai-ga}\\
    \gamma^\infty
    &=
    \frac{\upsilon\eta n}
    {1-\phi-\eta\upsilon}.
    \label{eq:cd-ai-gamma}
\end{align}
A positive constant-growth path requires
\begin{equation}
    1-\phi-\eta\upsilon>0.
    \label{eq:cd-ai-bg-condition}
\end{equation}
If this condition fails, the feedback from output to automated research and from
research back to output is too strong for a positive constant-growth solution. The
failure of \eqref{eq:cd-ai-bg-condition} identifies a candidate acceleration regime;
the transition system determines whether research and output accelerate or instead
move to a boundary allocation.

The remaining balanced-growth objects are
\begin{align}
    g_Y=g_K=g_C=g_U=g_M&=n+\gamma^\infty,
    \label{eq:aggregate-bg-rates}\\
    g_X&=n+\gamma^\infty+g_A^\infty,
    &
    g_w&=\gamma^\infty,
    &
    g_{p^X}&=-g_A^\infty,
    \label{eq:other-bg-rates}\\
    r^\infty&=\rho+\delta_K+\gamma^\infty,
    &
    \frac{K_t}{Y_t}
    &=
    \frac{\alpha}{\rho+\delta_K+\gamma^\infty}.
    \label{eq:bg-capital-output}
\end{align}
The costate equation and the research first-order condition imply
\begin{equation}
    m_R^\infty
    =
    \frac{\beta^2\eta g_A^\infty}
    {\rho-n+(1-\phi)g_A^\infty}.
    \label{eq:bg-research-resource-share}
\end{equation}
Consequently,
\begin{align}
    \frac{I_t}{Y_t}
    &=
    \frac{\alpha(n+\gamma^\infty+\delta_K)}
    {\rho+\delta_K+\gamma^\infty},
    \label{eq:bg-investment-share}\\
    \frac{C_t}{Y_t}
    &=
    1-
    \frac{\alpha(n+\gamma^\infty+\delta_K)}
    {\rho+\delta_K+\gamma^\infty}
    -\beta^2-m_R^\infty.
    \label{eq:bg-consumption-share}
\end{align}
An interior path requires a positive consumption share.

Because $g_w=\gamma^\infty>0$ and $\xi$ is constant,
Proposition~\ref{prop:gradual-automation} applies. More precisely,
\begin{equation}
    g_{H/N}
    =
    -(\theta-1)\gamma^\infty<0.
    \label{eq:human-research-population-share}
\end{equation}
The fraction of the population employed in frontier research therefore converges to
zero. The number of human researchers need not fall in levels: it grows at rate
$n-(\theta-1)\gamma^\infty$ and declines only if the second term exceeds population
growth.

\subsubsection{Why the production CES matters}

The Cobb--Douglas case is the knife-edge at which expenditure shares remain
constant. For $\sigma\neq1$, the monopoly condition
\eqref{eq:monopoly-ai-price} implies different asymptotic regimes as the marginal
cost $\xi/A_t$ falls.

\begin{proposition}[Production regimes under increasingly capable AI]
\label{prop:production-ces-regimes}
Suppose $A_t\to\infty$ and the service-market solution remains interior.
\begin{enumerate}[label=(\roman*)]
    \item If $\sigma<1$, the monopolist's AI share converges to
    \begin{equation}
        \bar s^X
        =
        \frac{1-\sigma}{1-\alpha\sigma}.
        \label{eq:low-sigma-ai-share}
    \end{equation}
    The ratio $X_t/L_{Y,t}$ remains bounded and the productive effect of further
    reductions in $\xi/A_t$ vanishes.
    \item If $\sigma=1$, factor shares are constant and the AI-dominated
    balanced-growth path is characterized by
    \eqref{eq:cd-ai-ga}--\eqref{eq:bg-consumption-share}.
    \item If $\sigma>1$, $s^X_t\to1$, the production-labor share converges to zero,
    and
    \begin{equation}
        \frac{Y_t}{K_t}
        \sim
        D_\sigma A_t^{(1-\alpha)/\alpha},
        \qquad D_\sigma>0.
        \label{eq:high-sigma-ak-limit}
    \end{equation}
    A conventional balanced-growth path with $g_A>0$ cannot have a constant
    capital--output ratio. The return to capital and the aggregate growth rate rise
    with capability.
\end{enumerate}
\end{proposition}

For $\sigma<1$, a rising wage relative to $\xi$ is not guaranteed: labor remains
essential and the monopoly supply of AI services approaches a finite
revenue-maximizing ratio. Research automation can therefore stop short of
$s^M=1$ even though $\theta>1$. For $\sigma=1$, the positive wage growth in
\eqref{eq:other-bg-rates} makes automated research asymptotically dominant. For
$\sigma>1$, production itself becomes an acceleration mechanism. Research
automation reinforces that mechanism, but the wage path must be solved jointly with
capital accumulation before concluding that the human research share disappears.

\subsubsection{Private and social research}

The planner's first-order condition for research compute is
\begin{equation}
    Q_tF_{M,t}=\xi.
    \label{eq:social-compute-foc}
\end{equation}
The private condition is $q_tF_{M,t}=\xi$. At the private allocation, the
uninternalized marginal social benefit is therefore
\begin{equation}
    Q_tF_{M,t}-\xi=(Q_t-q_t)F_{M,t}.
    \label{eq:knowledge-wedge}
\end{equation}

\begin{proposition}[The knowledge and appropriability wedge]
\label{prop:knowledge-wedge}
If $Q_t>q_t$, the integrated monopolist invests too little in research compute at the
margin relative to the social allocation.
\end{proposition}

There are thus two distinct distortions. Static market power restricts current AI
services, while incomplete appropriation of consumer surplus and knowledge benefits
can depress frontier research. Price regulation alone may correct the first
distortion while weakening the rents that finance research. This is the central
policy trade-off developed below.
